\section{Free source inputs and untouched exterior gates}

The separation argument in
\cite[proof of Theorem~5.2]{OpenAI2026PolynomialPEPS},
\texttt{04-compression.tex}, lines 342--434, uses corrected source
registers as free inputs. Exact sources internal to either side are
prepared there, while gates wholly in the exterior retain their
complete operators. We give the corresponding exact constructions.

% These entries provide the gate data, selective preparation, and owner
% identifications. The chronological partially expanded circuit and the
% corrected-subset trace-norm estimate are separate obligations.

\subsection{Finite gate expansions}

\begin{definition}[A gate with common finite source positions]
    \label{def:peps_prepared_source_gate}
    \lean{TNLean.PEPS.PairEffect.PreparedSourceGate}
    \lean{TNLean.PEPS.PairEffect.PreparedSourceGate.branchWord}
    \lean{TNLean.PEPS.PairEffect.PreparedSourceGate.eval}
    \leanok
    \uses{def:peps_source_slots,def:peps_party_registers}
    Fix a finite branch set $\Xi$, coefficients $c_\xi\in\C$, and
    input and output register lists $a,b$ on the gate's participating
    party set $P$. Common finite source data consist of an ordered list
    $E$ containing each unordered pair of distinct parties exactly once,
    spaces $U_e=\C^{u_e}$ and $V_e=\C^{v_e}$, normalized vectors
    $\eta_{\xi,e}\in U_e\otimes V_e$, and allowed compositions
    $T_\xi:\mathcal S\otimes\mathcal H_a\to\mathcal H_b$
    containing no pair preparations. Here $\mathcal S$ is the memory of
    all source registers in their prescribed order. Write
    $F_\xi=T_\xi Q_{\eta_\xi,\mathcal H_a}$ and
    \begin{align}
        G & =\sum_{\xi\in\Xi}c_\xi F_\xi,
        \label{eq:peps_prepared_gate_sum}
    \end{align}
    and require $\|G\|\leq1$. The coefficients are fixed before these
    data are chosen. Completion by unused pairs is taken only over the
    participating parties $P$ of this gate.
\end{definition}

\begin{theorem}[Prepared branches and the complete gate]
    \label{thm:peps_prepared_gate_properties}
    \lean{TNLean.PEPS.PairEffect.PreparedSourceGate.branchWord_isAllowed}
    \lean{TNLean.PEPS.PairEffect.PreparedSourceGate.sources_branchWord}
    \lean{TNLean.PEPS.PairEffect.PreparedSourceGate.norm_eval_le_one}
    \leanok
    \uses{def:peps_prepared_source_gate}
    Each $F_\xi$ is an allowed composition. Its ordered source list is
    precisely the list of prescribed pairs with spaces $U_e,V_e$ and
    vectors $\eta_{\xi,e}$. The complete sum $G$, with its original
    scalar coefficients, is a contraction.
\end{theorem}

\begin{proof}
    \leanok
    \uses{thm:peps_source_preparation,def:peps_prepared_source_gate}
    Normalization makes the source preparation allowed. Compose it with
    $T_\xi$, whose source list is empty. Thus the only sources are the
    prescribed preparations. The bound on $G$ is part of the gate data.
\end{proof}

\begin{theorem}[Construction from the original branch operators]
    \label{thm:peps_prepared_gate_exists}
    \lean{TNLean.PEPS.PairEffect.Word.exists_preparedSourceGate}
    \leanok
    \uses{def:peps_prepared_source_gate}
    Suppose $P$ and $\Xi$ are finite, each $W_\xi:a\to b$ is an
    allowed composition, and
    $\|\sum_{\xi\in\Xi}c_\xi W_\xi\|\leq1$.
    There are common finite source data with the same coefficients
    $c_\xi$ for which $G=\sum_{\xi\in\Xi}c_\xi W_\xi$ exactly.
    No nonemptiness assumption on $\Xi$ is needed.
\end{theorem}

\begin{proof}
    \leanok
    \uses{thm:peps_finite_source_gate}
    Apply Theorem~\ref{thm:peps_finite_source_gate} to the original
    branch compositions. It gives common finite spaces, normalized
    vectors, and allowed compositions without sources, with
    $T_\xi Q_{\eta_\xi,\mathcal H_a}=W_\xi$ for each $\xi$.
    Substitution in (\ref{eq:peps_prepared_gate_sum}) preserves the
    original sum and therefore its contraction bound.
\end{proof}

\subsection{Affected parties and the exterior}

\begin{definition}[Keeping affected owners distinct]
    \label{def:peps_affected_owner}
    \lean{TNLean.PEPS.PairEffect.affectedOwner}
    \lean{TNLean.PEPS.PairEffect.affectedOwner_eq_none}
    \leanok
    Let $\mathcal A\subseteq P$ have decidable membership. Set
    $Q=\mathcal A\sqcup\{*\}$, where $*$ denotes the exterior, and
    define $g:P\to Q$ by
    \begin{align}
        g(p) & =
        \begin{cases}
            p, & p\in\mathcal A,    \\
            *, & p\notin\mathcal A.
        \end{cases}
        \label{eq:peps_affected_owner}
    \end{align}
    Thus $g(p)=*$ exactly when $p$ is exterior.
\end{definition}

\begin{theorem}[Equality of grouped owners]
    \label{thm:peps_affected_owner_equality}
    \lean{TNLean.PEPS.PairEffect.affectedOwner_eq_iff}
    \lean{TNLean.PEPS.PairEffect.PairSource.affectedOwner_ne_iff}
    \leanok
    \uses{def:peps_affected_owner}
    One has $g(p)=g(q)$ exactly when $p=q$ or both parties are
    exterior. Consequently, for distinct source endpoints $p,q$,
    $g(p)\neq g(q)$ exactly when at least one endpoint belongs to
    $\mathcal A$. Both corrected pairs within $\mathcal A$ and pairs
    crossing to its exterior remain between distinct owners.
\end{theorem}

\begin{proof}
    \leanok
    \uses{def:peps_affected_owner}
    Separate the four possibilities for membership of $p,q$ in
    $\mathcal A$. Two affected images agree only when their original
    parties agree; an affected image cannot equal $*$.
\end{proof}

\begin{theorem}[The surviving ordered source occurrences]
    \label{thm:peps_affected_source_occurrences}
    \lean{TNLean.PEPS.PairEffect.SourceInventory.mem_mapOwner_affectedOwner}
    \lean{TNLean.PEPS.PairEffect.SourceInventory.map_partyPair_affectedOwner}
    \lean{TNLean.PEPS.PairEffect.SourceInventory.mapOwner_affectedOwner_eq_nil_iff}
    \leanok
    \uses{def:peps_affected_owner,def:peps_owner_sources}
    A source occurs in $S_g$ precisely when it is obtained from an
    occurrence $s=(p,q;U,V;\eta)$ in $S$ with at least one affected
    endpoint, by replacing its owners with $g(p),g(q)$. Its spaces and
    vector are unchanged. The ordered pair labels are obtained by
    filtering $S$ by this same condition and then applying $g$ to both
    endpoints. Order and multiplicity are retained, even when different
    original pairs acquire the same image. In particular, $S_g$ is
    empty exactly when both endpoints of every original source are
    exterior.
\end{theorem}

\begin{proof}
    \leanok
    \uses{def:peps_owner_sources,thm:peps_owner_sources_internal,
        thm:peps_affected_owner_equality}
    Substitute the distinct-owner criterion of
    Theorem~\ref{thm:peps_affected_owner_equality} into the occurrence
    description of $S_g$. Induction on the original source list gives
    the ordered pair-label identity. The empty-list criterion follows
    from Theorem~\ref{thm:peps_owner_sources_internal}.
\end{proof}

\begin{definition}[Retaining a complete operation at one grouped owner]
    \label{def:peps_grouped_block_map}
    \lean{TNLean.PEPS.PairEffect.Word.groupedBlockMap}
    \leanok
    \uses{def:peps_owner_layout,def:peps_party_layout}
    Let $f:P\to Q$, and suppose every register in $a$ and $b$ has
    owner image $q\in Q$. For $D:\mathcal H_a\to\mathcal H_b$,
    regard $R_{f,b}DR_{f,a}^{-1}$ as a local operation at $q$ and
    leave the spectator list $\ell$ untouched. This gives an actual
    composition $D^f_{\ell}:(a\ell)^f\to(b\ell)^f$.
\end{definition}

\begin{theorem}[Exact action of an untouched aggregate]
    \label{thm:peps_grouped_block_map}
    \lean{TNLean.PEPS.PairEffect.Word.groupedBlockMap_spec}
    \lean{TNLean.PEPS.PairEffect.Word.eval_groupedBlockMap}
    \leanok
    \uses{def:peps_grouped_block_map}
    The complete operation satisfies
    \begin{align}
        D^f_{\ell}R_{f,a\ell}
         & =R_{f,b\ell}C_{b,\ell}^{-1}
        (D\otimes I_{\mathcal H_\ell})C_{a,\ell}.
        \label{eq:peps_grouped_block_action}
    \end{align}
    If $\|D\|\leq1$, this composition is allowed and contains no pair
    sources. For $f=g$ and $q=*$, this retains an exterior gate as its
    original aggregate contraction, without expanding its branch sum.
\end{theorem}

\begin{proof}
    \leanok
    \uses{thm:peps_owner_concatenation,def:peps_grouped_block_map}
    Apply the concatenation identity for the canonical owner
    isometries. The two isometries conjugating $D$ cancel with the
    corresponding input and output identifications, giving
    (\ref{eq:peps_grouped_block_action}). Conjugation by isometries
    preserves the contraction bound. A local operation creates no
    pair sources.
\end{proof}

\begin{theorem}[Grouping preserves the weighted gate sum]
    \label{thm:peps_grouped_gate_sum}
    \lean{TNLean.PEPS.PairEffect.Word.eval_sum_mapOwner_comp}
    \leanok
    \uses{def:peps_owner_composition}
    For a finite family of compositions $F_\xi:a\to b$ and arbitrary
    coefficients $c_\xi\in\C$,
    \begin{align}
        \left(\sum_\xi c_\xi F_\xi^f\right)R_{f,a}
         & =R_{f,b}\left(\sum_\xi c_\xi F_\xi\right).
        \label{eq:peps_grouped_gate_sum}
    \end{align}
\end{theorem}

\begin{proof}
    \leanok
    \uses{thm:peps_owner_composition}
    Sum the exact identity $F_\xi^fR_{f,a}=R_{f,b}F_\xi$ with
    coefficients $c_\xi$. Composition is linear on both sides.
\end{proof}

\begin{definition}[The complete gate after grouping]
    \label{def:peps_prepared_gate_grouped}
    \lean{TNLean.PEPS.PairEffect.PreparedSourceGate.evalAtOwners}
    \leanok
    \uses{def:peps_prepared_source_gate,def:peps_owner_layout}
    For common finite source data and any map $f:P\to Q$, define
    the grouped aggregate by $G^f=R_{f,b}GR_{f,a}^{-1}$.
\end{definition}

\begin{theorem}[The original expansion and aggregate norm after grouping]
    \label{thm:peps_prepared_gate_grouped}
    \lean{TNLean.PEPS.PairEffect.PreparedSourceGate.evalAtOwners_eq_sum}
    \lean{TNLean.PEPS.PairEffect.PreparedSourceGate.norm_evalAtOwners_le_one}
    \leanok
    \uses{def:peps_prepared_gate_grouped,def:peps_owner_composition}
    The exact branch expansion is $G^f=\sum_\xi c_\xi F_\xi^f$,
    and $\|G^f\|\leq1$. This bound comes from the original
    aggregate contraction and does not involve
    $\sum_\xi|c_\xi|$.
\end{theorem}

\begin{proof}
    \leanok
    \uses{thm:peps_grouped_gate_sum,thm:peps_prepared_gate_properties}
    Compose (\ref{eq:peps_grouped_gate_sum}) with $R_{f,a}^{-1}$.
    Both canonical owner maps are isometries, so conjugating the
    aggregate preserves its norm bound.
\end{proof}

\subsection{Preparing only the fixed sources}

\begin{definition}[The selected original source positions]
    \label{def:peps_selected_source_positions}
    \lean{TNLean.PEPS.PairEffect.SourceInventory.selectedSources}
    \lean{TNLean.PEPS.PairEffect.SourceInventory.layout_selectedSources_eq}
    \lean{TNLean.PEPS.PairEffect.SourceInventory.freeSlotLayout}
    \lean{TNLean.PEPS.PairEffect.SourceInventory.mem_selectedSources}
    \leanok
    \uses{def:peps_source_slots}
    Let $E$ be an ordered list of pair positions, with halfspaces
    $U_e,V_e$, and let $E_{\rm free}\subseteq E$ be selected. The
    selected source list retains these occurrences in their original
    order, with their original endpoint parties and the prescribed
    spaces and vectors. Its register layout is independent of those
    vectors. Denote its memory by $\mathcal S_{\rm free}$ and the
    memory of all positions by $\mathcal S$.
\end{definition}

\begin{definition}[Selective source preparation]
    \label{def:peps_selective_preparation}
    \lean{TNLean.PEPS.PairEffect.SourceInventory.prepareSelected}
    \lean{TNLean.PEPS.PairEffect.SourceInventory.prepareFreeSlots}
    \lean{TNLean.PEPS.PairEffect.SourceInventory.fillSourceVectors}
    \leanok
    \uses{def:peps_selected_source_positions}
    Fix vectors $\eta_e\in U_e\otimes V_e$ only at positions outside
    $E_{\rm free}$. Define
    $\Phi_\eta:\mathcal S_{\rm free}\otimes\mathcal H_\ell
        \to\mathcal S\otimes\mathcal H_\ell$ by preparing the fixed
    sources and retaining all free registers in their original
    positions. This is one actual composition, independent of the
    vectors later supplied to the free registers.

    For an arbitrary family $\zeta_e\in U_e\otimes V_e$, write
    $Q^{\rm free}_{\zeta,\mathcal H_\ell}$ for preparation at just
    the free positions. The completed family $\theta$ has
    $\theta_e=\zeta_e$ at free positions and $\theta_e=\eta_e$
    elsewhere.
\end{definition}

\begin{theorem}[Only the fixed sources are prepared]
    \label{thm:peps_selective_preparation_sources}
    \lean{TNLean.PEPS.PairEffect.SourceInventory.isAllowed_prepareSelected}
    \lean{TNLean.PEPS.PairEffect.SourceInventory.sources_prepareSelected}
    \lean{TNLean.PEPS.PairEffect.SourceInventory.mem_sources_prepareSelected}
    \leanok
    \uses{def:peps_selective_preparation}
    The source list of $\Phi_\eta$ is exactly the list of fixed
    original positions, in their original order, with vectors
    $\eta_e$. In particular, each prepared source has the original
    endpoint parties and halfspaces of its fixed position. If every
    fixed vector has norm one, then $\Phi_\eta$ is allowed.
\end{theorem}

\begin{proof}
    \leanok
    \uses{def:peps_selective_preparation,def:peps_selected_source_positions}
    Induct on the source list. At a free position, retain the two
    registers and apply the construction to the remaining positions.
    At a fixed position, prepare its source after the remaining
    positions. The source inventory records precisely this fixed
    occurrence. Its preparation is allowed by normalization.
\end{proof}

\begin{theorem}[Uniform recovery of every completed preparation]
    \label{thm:peps_selective_preparation_recovery}
    \lean{TNLean.PEPS.PairEffect.SourceInventory.eval_prepareSelected_prepareFreeSlots}
    \leanok
    \uses{def:peps_selective_preparation}
    For every family of free vectors, without a normalization
    assumption,
    \begin{align}
        \Phi_\eta Q^{\rm free}_{\zeta,\mathcal H_\ell}
         & =Q_{\theta,\mathcal H_\ell}.
        \label{eq:peps_selective_recovery}
    \end{align}
\end{theorem}

\begin{proof}
    \leanok
    \uses{def:peps_selective_preparation,thm:peps_source_preparation}
    Induct on the ordered positions. At a free head, preparation of
    its vector commutes with the operations on the remaining
    registers. At a fixed head, both sides prepare the same fixed
    vector. In each case the induction hypothesis gives the equality
    for the remaining positions. The empty source list gives the
    identity on the original input memory.
\end{proof}

\begin{theorem}[Two contractions independent of the free source vectors]
    \label{thm:peps_selective_partition}
    \lean{TNLean.PEPS.PairEffect.Word.exists_selective_partition}
    \leanok
    \uses{def:peps_selective_preparation,def:peps_party_partition}
    Let $f:P\to\{0,1\}$ specify two sides. Suppose every fixed
    source has both endpoints on the same side and has norm one.
    Let $T:\mathcal S\otimes\mathcal H_\ell\to\mathcal H_{\ell'}$
    be an allowed composition containing no pair sources. There are
    two allowed compositions $A_+,A_-$, containing no pair sources and
    each having norm at most one, such that
    \begin{align}
        J_{f,\ell'}T\Phi_\eta
         & =(A_+\otimes A_-)J_{f,L},
        \label{eq:peps_selective_partition} \\
        J_{f,\ell'}TQ_{\theta,\mathcal H_\ell}
         & =(A_+\otimes A_-)J_{f,L}
        Q^{\rm free}_{\zeta,\mathcal H_\ell}.
        \label{eq:peps_selective_all_vectors}
    \end{align}
    Here $L$ is the free-source register list followed by $\ell$, and
    $J_{f,L}$ is its canonical partition isometry. The same two
    compositions work for every family $\zeta$ of free vectors.
    Thus the contraction bounds apply on all free inputs.
\end{theorem}

\begin{proof}
    \leanok
    \uses{thm:peps_selective_preparation_sources,
        thm:peps_internal_source_factorization,
        thm:peps_selective_preparation_recovery}
    The only sources in $T\Phi_\eta$ are the fixed original
    positions. They are internal by hypothesis. Apply
    Theorem~\ref{thm:peps_internal_source_factorization} to obtain
    (\ref{eq:peps_selective_partition}) and the two contraction
    bounds. Compose that identity with
    $Q^{\rm free}_{\zeta,\mathcal H_\ell}$ and use
    (\ref{eq:peps_selective_recovery}) to obtain
    (\ref{eq:peps_selective_all_vectors}). The factors were chosen
    before $\zeta$.
\end{proof}

\subsection{Operations beside untouched registers}

\begin{definition}[Extension by spectator registers]
    \label{def:peps_append_spectators}
    \lean{TNLean.PEPS.PairEffect.Word.appendTail}
    \leanok
    \uses{def:peps_party_layout,thm:peps_register_exchange}
    Given an actual composition $W:a\to b$ and a spectator list
    $\ell$, move $\ell$ before $a$, perform $W$ while leaving those
    registers unchanged, and move them back after $b$. Denote the
    resulting composition $a\ell\to b\ell$ by $W_{\ell}$.
\end{definition}

\begin{theorem}[Exact spectator extension]
    \label{thm:peps_append_spectators}
    \lean{TNLean.PEPS.PairEffect.Word.isAllowed_appendTail}
    \lean{TNLean.PEPS.PairEffect.Word.sources_appendTail}
    \lean{TNLean.PEPS.PairEffect.Word.eval_appendTail}
    \leanok
    \uses{def:peps_append_spectators}
    The spectator extension has exactly the same ordered source list
    as $W$, is allowed whenever $W$ is allowed, and satisfies
    \begin{align}
        W_{\ell}C_{a,\ell}^{-1}
         & =C_{b,\ell}^{-1}(W\otimes I_{\mathcal H_\ell}).
        \label{eq:peps_append_spectators}
    \end{align}
    The operator identity holds without norm or normalization
    assumptions.
\end{theorem}

\begin{proof}
    \leanok
    \uses{thm:peps_register_exchange,thm:peps_source_preparation}
    On an elementary input $x\otimes y$, the first exchange gives
    $y\otimes x$, the framed composition gives $y\otimes Wx$, and
    the last exchange gives $Wx\otimes y$. Linearity gives
    (\ref{eq:peps_append_spectators}). Exchanges create no sources,
    and leaving a register unchanged preserves the source list and
    allowedness of the operations on the other registers.
\end{proof}
